% DUET -- Introduction (DAC 2027 draft, 2026-10-20)
% Numbers: results/paper_dac/tables.md (NVMeVirt, 12.5% mapping DRAM, OP 7%). TODO-marked cites need sources.

\section{Introduction}
\label{sec:intro}

The flash translation layer (FTL) of an SSD maps every logical page to its
physical location. A page-level mapping table with 4\,B entries for 4\,KiB pages
takes about 0.1\% of the capacity, so a 32\,TB SSD needs 32\,GB of DRAM for
the table alone. SSD capacity keeps growing with denser 3D NAND, but controller
DRAM does not grow with it because of its cost, power, and board area~\cite{TODO-dram-cost}.
Large SSDs therefore have to translate addresses with a mapping DRAM budget
that is only a fraction of the full page-level table.

Classic FTLs fit a limited budget by keeping only part of the mapping in DRAM,
and pay for it in latency.
\emph{On-demand mapping} (DFTL~\cite{dftl}, LeaFTL~\cite{leaftl}) caches the
page-level table and fetches missing entries from flash before the data read.
With 1/8 of the table's DRAM, LeaFTL's p99 read latency on YCSB-A is
3.7$\times$ that of a page-level FTL with unlimited DRAM.
\emph{Hybrid mapping} (FAST~\cite{fast}) maps data per erase block and merges
page-mapped log blocks back into it; a merge copies up to a whole block, and
blocks keep growing with each 3D NAND generation~\cite{TODO-block-size}.
FAST's p99 write latency is 2.8--4.7$\times$ that of a page-level FTL under random I/O.

To avoid these latency problems, high-capacity SSDs in production use a
\emph{large indirection unit} (IU) of 16--64\,KiB instead~\cite{TODO-iu}: the
whole coarse table fits in DRAM, with no mapping misses and no block merges.
The cost moves to small writes, each of which becomes a read-modify-write (RMW)
of its unit; on YCSB this raises the write amplification factor (WAF) by 2.7--3.3$\times$.

We present \textbf{DUET} (\textbf{DU}al-granularity \textbf{E}lastic
\textbf{T}ranslation), an FTL that keeps the large IU but removes its RMW
by mapping data at two granularities under a fixed DRAM budget. Like hybrid
mapping it has a coarse and a fine level, but its coarse unit is the 64\,KiB IU,
decoupled from the erase block. DUET has three key elements:
\begin{itemize}
  \item \textbf{Dual-granularity mapping.} Most data are mapped per 64\,KiB IU,
  which keeps the coarse table at 1/16 of a page-level table. Writes smaller
  than an IU are written out-of-place as 4\,KiB fragments, which a compact cuckoo
  hash index in DRAM tracks, so there is no RMW on the write path.
  \item \textbf{Bounded, cost-aware folding.} Fragments are merged into their IU
  only when the merge pays off, and most are merged for free while garbage
  collection relocates the IU anyway. A merge touches at most one IU (16 pages),
  regardless of the erase-block size, and space is reclaimed by page-style
  garbage collection rather than block merges, so no merge ever stalls host
  writes for the length of a block copy.
  \item \textbf{Elastic fragment mapping.} When the DRAM index is full, DUET
  spills fragment mappings to flash instead of forcing merges, and brings them
  back into DRAM when the IU is accessed again. Read-promoted mappings stay
  \emph{clean}: their flash copy remains valid, so evicting them again costs no flash write.
  The fine-grained mapping thus stretches beyond DRAM and shrinks back without
  data movement.
\end{itemize}

We implement DUET and six other FTLs in NVMeVirt~\cite{nvmevirt}, a kernel-level
NVMe SSD emulator, with the same NAND timing, write buffer, and garbage collection
for all of them. We evaluate them with a 1/8 mapping-DRAM budget on real traces
(Alibaba cloud block storage~\cite{alibaba-trace} and SPFresh~\cite{spfresh}), key-value and OLTP
applications (YCSB~\cite{ycsb} on PostgreSQL and TPC-C), and raw random I/O.
DUET reaches 92\% of the throughput of a page-level FTL with unlimited DRAM (geometric mean),
compared to 79\% for LeaFTL and 68\% for DFTL. On the Alibaba trace it is
18\% faster than the page-level FTL. On the read side, its average read latency is up to 5.5$\times$ and
2.5$\times$ lower than DFTL's and LeaFTL's, respectively; on the write side, its p99 write latency is
up to 4.1$\times$ lower than FAST's (1.2--4.1$\times$ on every workload with a meaningful write share).

This paper makes the following contributions:
\begin{itemize}
  \item We show that existing FTLs fit a limited mapping DRAM budget only by
  paying in read latency (mapping caching), merge cost and tail latency
  (block-level hybrids), or RMW (coarse IUs).
  \item We design DUET, which combines dual-granularity mapping, bounded cost-aware
  folding, and elastic spilling and clean promotion of fragment mappings.
  \item We implement seven FTLs in a common emulator and show that DUET outperforms
  all of them under the same DRAM budget.
\end{itemize}
